\section{Exact arithmetic and reproduction of the grid certificate}
\label{grid:implementation}

The complete C++17 program in Section~\ref{cert:grid} implements the
recursions of Section~\ref{grid:section}.  Its parameters are fixed at
$H=72$, $N=10^{12}$, and $\mathtt{SCALE}=10^{15}$.  It uses signed
64-bit integers for table coefficients and mass units, and signed
128-bit integers for products and accumulators. The compiler must also
provide a plain \texttt{int} type of at least 32 bits. Assertions must be
enabled; in particular the build must not define \texttt{NDEBUG}.
For example, from the paper directory, the following commands compile and run all 28 cases:
\begin{verbatim}
g++ -std=c++17 -O3 -UNDEBUG -Wall -Wextra \
    verification/grid.cpp -o grid
for j in $(seq -1 26); do ./grid "$j"; done
\end{verbatim}
The program writes progress diagnostics to standard error and the case
number followed by its stopping iteration to standard output.  Only
the diagnostics convert to floating point.  Successful outputs are
specified in~\eqref{grid:outputs}, together with \texttt{-1 162}.

\subsection{Correspondence of code and probability operations}

The routine \texttt{init} enumerates all sorted nonnegative integer
triples of sum $72$.  It finds $469$ triples, identifies the pure
representative $(72,0,0)$ and the uniform representative $(24,24,24)$,
and stores an index for every representative.  The routine
\texttt{indexof} sorts a candidate triple and verifies the resulting
index.  A \texttt{Law} is a vector of orbit masses in units of $1/N$.
The routine \texttt{complete} adds the missing units at the pure orbit,
which is exactly Lemma~\ref{grid:completion}.

The constructor \texttt{Table(mode,A,AD)} first reduces the channel
fraction.  The variables \texttt{lhs}, \texttt{rhs}, \texttt{raw},
\texttt{T}, and \texttt{den} are the entries and denominator of
\eqref{grid:modes}.  For each of the six relative permutations,
\texttt{n}, \texttt{rem}, and \texttt{miss} are the floor vector,
remainders, and deficit in the grid spread.  The corner weights are
\texttt{T}, \texttt{rem[c]}, or \texttt{T-rem[c]}, as appropriate.
After aggregating equal output indices, the constructor checks that
their sum is \texttt{den}.  These checks include every grid pair, not
only pairs visited by an iteration.  Zero normalizers are omitted
because they have zero tilted mass.

The routine \texttt{combine} implements~\eqref{grid:combine}.  Symmetric
modes use the triangular pair table and the stated off-diagonal
coefficient; mode~1 uses ordered pairs.  Its row-skipping condition is
safe: an ordered row contributes zero when its first mass is zero,
while a triangular row contributes zero when both first-index masses
are zero.  The table position advances by the full skipped row length.
After division by $N\,\mathtt{den}$, each mass is rounded down and
\texttt{complete} is applied.

The functions \texttt{down} and \texttt{up} compute floor and ceiling
of a nonnegative integer ratio.  The latter uses
$\lceil x/y\rceil=\lfloor(x+y-1)/y\rfloor$ for $x\geq0$, $y>0$.
The function \texttt{w\_zero} implements~\eqref{grid:taylor}.
The routine \texttt{poisson} fixes its incoming child law
\texttt{inc}, starts \texttt{part} at the uniform zero-count law, and
uses mode~1 once for each positive count weight.  Its orbit sums are
exactly~\eqref{grid:mixture}; the final \texttt{complete} handles all
omitted mass.  No count product is needed after the first zero weight.
The regular branch of \texttt{run} is~\eqref{grid:regular-recursion}.
Finally, \texttt{mu\_num} is $M$ in~\eqref{grid:moment}, and the strict
stopping test is $20M<NH^2$.  Each run allows $400$ iterations and
asserts failure if none succeeds.  All reported stopping iterations
are below that bound.

\subsection{Bounds on every arithmetic class}

For all Poisson table calls, the reduced fraction satisfies
$B\leq10^6$ and $A\leq B/2$.  Thus
\[
 0\leq e_i\leq(B+A)H\leq108000000,
 \quad D_{\mathrm{tab}}=6BH^2\leq31104000000,
 \quad Hr_i\leq559872000000.
\]
Each total $T$ is at most three times the bound on a raw coordinate;
all of these quantities fit signed 64 bits.  The only mode~2 call is
at the reduced fraction $A/B=1/2$, where
$D_{\mathrm{tab}}=373248$ and $Hr_i\leq3359232$.
Mode~0 has smaller bounds.  Every nonnegative aggregated table weight
and every partial sum of weights is bounded by its denominator, by
\eqref{grid:table-mass}.

Input masses are nonnegative and sum to $N$, so their ordered pair
mass is $N^2$.  The triangular implementation has the same total
after the off-diagonal terms are combined.  Consequently every
nonnegative product accumulator and its partial sums are at most
\[
 D_{\mathrm{tab}}N^2
 \leq31104000000000000000000000000000000<2^{127}-1.
\]
The source casts an input factor to the 128-bit type \texttt{Big}
before multiplying.  Dividing by $ND_{\mathrm{tab}}$ yields a mass
between zero and $N$, safely convertible back to 64 bits.  The
uncompleted masses sum to at most $N$; likewise the Poisson mixture
masses sum to at most $N$, since they are lower bounds on a probability
mixture.  These facts bound all 64-bit partial sums in
\texttt{complete} and \texttt{poisson}.

Here are bounds for the Taylor calculation independent of machine
floating point.  If $t_n=S_0d_+^n/n!$, the lower and upper endpoint
errors satisfy
\[
 0\leq t_n-l_n< e^{d_+},\qquad
 0\leq u_n-t_n< e^{d_+}.
\]
Indeed an error satisfies $\epsilon_n\leq(d_+/n)\epsilon_{n-1}+1$;
expanding this recurrence bounds it by
$\sum_{j\geq0}d_+^j/j!=e^{d_+}$, since each product of $j$ successive
denominators is at least $j!$.  As $d_+<5$ and $e<3$, the endpoint
errors are less than $243$.  The sum of all endpoint magnitudes through
degree $41$ is therefore bounded by
\[
 S_0e^{d_+}+42\cdot243<243(S_0+42)<2^{63}-1.
\]
This also bounds every signed partial sum in~\eqref{grid:taylor}.
The endpoint multiplication by $d_{\max}\leq881$ is carried out in
128 bits before division, and the same holds for multiplication by
$N$ and for the Poisson weight recurrence.  These products are all
less than $2^{127}-1$ under the displayed bounds.

The lower sum $L$ is positive.  In fact Taylor's theorem gives
\[
 \frac{L}{S_0}
 \geq e^{-d_+}-\frac{d_+^{42}}{42!}
                   -\frac{41\cdot243}{S_0}
 >\frac1{243}-10^{-20}-\frac{9963}{10^{15}}>0.
\]
The middle inequality uses the elementary integer comparison
$5^{42}/42!<10^{-20}$.  Thus the nonnegative-argument conditions of
\texttt{down} hold at the final Taylor step as well as in its term
recurrences. Since $w_n/N\leq p_n$, every Poisson weight is at most
$N$. In particular $w_8\leq N<2^{40}$, and the recurrence gives
$w_n\leq w_{n-1}/2$ for $n\geq9$, since $d_+/n<1/2$.
Thus $w_{48}<1$, so the first zero occurs by $n=48$. The count loop
and its denominator therefore satisfy $n\leq48$ and $200n\leq9600$,
well within the plain integer type. Every product $w_nq_n(h)$ is at
most $N^2$ and is evaluated in 128 bits.

For $h\in\mathcal H$, the bracket in~\eqref{grid:moment} belongs to
$[0,H^2]$.  Hence every partial sum in \texttt{mu\_num} is at most
\[
 NH^2=5184000000000000<2^{63}-1.
\]
The factor $20$ in the stopping comparison is applied after converting
to \texttt{Big}.  The binary-search comparisons for~\eqref{grid:magnitude}
likewise use 128 bits before taking integer squares and products.

There are at most $18$ unaggregated contributions per grid pair
(six permutations, at most three corners).  Thus even a full ordered
table has at most $18\cdot469^2=3959298$ entries before aggregation.
All table positions and vector indices fit an ordinary signed 32-bit
integer, and the source also checks an explicit
$100000000$ entry limit.  Candidate grid coordinates belong to
$\{0,\ldots,H\}$, and all small grid sums and products fit these
integer types.  This accounts for the unguarded indexing arithmetic
as well as the probabilistic accumulators.

The supplied source therefore implements a finite exact-arithmetic
certificate: positivity, normalization, rounding direction, all
parameter choices, and the final strict inequalities are independent
of compiler floating-point behavior.
